% Successor calculation. Earlier sealed editions remain unchanged.
\begingroup
\section{Original conductor pivots and correlated directional receivers}
\label{sec:CRP}
\subsection{Objects, measures, and the attained fibre}
This calculation integrates the first corrected-row note, equations
(1)--(25), in the supplied transcript \texttt{Untitled 1775.md}, lines
1--851. The coefficient estimates are the fully proved LT1--22/LRC6--9
of \path{LRC_TRANSPORT_COMPLETE.tex}; the original affine fibre is
ATA1--8 of \path{ATTAINED_FIBRE_AND_PROFILE_PROPAGATION.tex}.
All programme and source attributions in those providers are retained.
Here every new step connecting those estimates to individual attained
pivots and their correlated directional sums is proved explicitly.

Fix the actual simple quartet, period, logarithm branch and conductor.
Write
\[
\begin{gathered}
a\equiv1\pmod4,\quad a\ge257,\quad 0<\delta<\tfrac12,\quad\gamma>2,\\
q=(a+1)^2,\quad Q=(a-7)^2,\quad\Delta=16a-48,\\
\beta_{b,u,w}=b/2+(2u-b)\delta+i(2w-b)\gamma,
\quad\chi_b(S)=\prod_{u,w=0}^{b}(S-\beta_{b,u,w}),\\
E_A(z)=\sum_{u,w=0}^{8}a_{uw}e^{\beta_{8,u,w}z},\quad
\mu_j=E_A^{(j)}(0),\quad v=\operatorname{ord}_0 E_A\le80,
\quad\mu_v\ne0,\\
s\in\{1,a\},\quad c=a/2,\quad c'=a/2-4,\quad
m=\Delta-v,\quad L=q-v,\\
0\le r\le M\le q+32a,\quad N=q+M,\quad D=N-v,\quad d=D+1.
\end{gathered}\tag{CRP1}
\]
The finite analytic root guards of LRC13 are imposed on both actual
grids, with the same original order $s$: for $b=a,a-8$,
\[
n_b=(b+1)^2+(s-1)/4\ge100,\qquad
\max\{b\sqrt{\delta^2+\gamma^2},(s-1)/2+1/2\}\le2^{-17}n_b.
\tag{CRP2}
\]
These are the domain of the supplied finite single-norm estimates;
the exact matrix arguments below require no additional root guard.
The original Hilbert space is
\[
\|P\|_{s,e}^2=\int_{\mathbb R}|P(e+iy)|^2dm_s(y),\qquad
dm_s(y)=\frac{M_s2^{s/2}}{4\pi\Gamma(s/2)}
 |\Gamma(s/4+iy/2)|^2dy,\quad M_s=(2\pi)^{s/2}.\tag{CRP3}
\]
No half-line, mass or centre has been changed. The original orthonormal
polynomials, including their phases, are
\[
\begin{gathered}
p_{0,s}=1,\ p_{1,s}=y,\quad
p_{n+1,s}=yp_{n,s}-n(n-1+s/2)p_{n-1,s},\\
h_{n,s}=M_sn!(s/2)_n,\quad
\phi_{n,s,e}(S)=p_{n,s}((S-e)/i)/\sqrt{h_{n,s}},\quad
\rho_n=\sqrt{n!/(s/2)_n},\\
\sum_{n\ge0}\frac{p_{n,s}(y)}{n!}t^n
 =(1+t^2)^{-s/4}e^{y\arctan t}.\tag{CRP4}
\end{gathered}
\]
The generating and orthogonality formulas are the original
Meixner--Pollaczek dictionary in NIST DLMF (18.23.7), (18.19.8)--(18.19.9)
\cite{human:dlmf18}, with $p_n(y)=n!P_n^{(s/4)}(y/2;\pi/2)$.
LT2--3 proves the factors in this dictionary, including $dy=2dx$ and
the original mass $M_s$.

The literal conductor is $(\mathcal T_AP)(S')=
\sum a_{uw}P(S'+\beta_{8,u,w})$. Since
$\mathcal T_AS^n=\sum_{j=v}^n\binom nj\mu_j S'^{n-j}$, its full
kernel on $\mathbb C[S]_{\le N}$ is $\mathbb C[S]_{<v}$ and it is
onto $\mathbb C[S']_{\le D}$. In the original quotient frame
$\phi_{v,s,c},\ldots,\phi_{N,s,c}$ and target frame
$\phi_{0,s,c'},\ldots,\phi_{D,s,c'}$, its square matrix is
\[
\begin{gathered}
A=A_{D,s},\qquad A_{ij}=\rho_i^{-1}g_{j-i}\rho_{j+v}\quad(i\le j),
\quad A_{ij}=0\quad(i>j),\\
w(t)=-i\arctan t,\quad g(t)=t^{-v}e^{-4w(t)}E_A(w(t)),
\quad g_0=(-i)^v\mu_v/v!.\tag{CRP5}
\end{gathered}
\]
This is LT4--6, obtained by applying the conductor to CRP4 and comparing
coefficients. The phase of $g_0$, all zeros of $g$, and all actual
coefficients remain in this matrix.

Let $P_r^+$ minimize $\|\chi_aP\|_{s,c}^2$ over monic degree-$r$
polynomials, and put $\nu_r^+=\|\chi_aP_r^+\|_{s,c}^2$.
Existence, uniqueness and orthogonality follow by completing the
positive coefficient square: the minimizer is orthogonal to
$\chi_a\mathbb C[S]_{<r}$. Consequently the original coordinate
columns
\[
\begin{gathered}
\Phi_r=\chi_aP_r^+/\sqrt{\nu_r^+},\quad
\Phi=(\Phi_0,\ldots,\Phi_M),\quad\Phi^*\Phi=I_{M+1},\\
\quad E=[0_{d\times v},I_d],\quad B=E\Phi,
\quad K=B^*B,\quad H=B^*A^*AB
\end{gathered}\tag{CRP6}
\]
are well defined in the upper native degree-$N$ frame. The matrix $E$
is precisely the orthogonal removal of the complete conductor kernel.
The map $B$ is injective: any element of $\chi_a\mathbb C[S]_{\le M}$
in $\mathbb C[S]_{<v}$ is zero, because $q>v$. Thus $0<K\preceq I$.

Write $u_r=\mathcal T_A(\chi_aS^r)/\ell_r$ with
$\ell_r=\mu_v\binom{q+r}{v}$, and retain the complete minima
\[
\eta_r=\min_{z\in\mathbb C^r}\|u_r+\sum_{t<r}z_tu_t\|_{s,c'}^2,
\quad\nu^-_{m+r}=\min_{P\ {
m monic},\deg P=m+r}
 \|\chi_{a-8}P\|_{s,c'}^2,
\quad\gamma_r=\log(\eta_r/\nu^-_{m+r})\ge0.\tag{CRP7}
\]
The original relation polynomials lie in this lower ideal. For
$P=S^r+\sum_{t<r}p_tS^t$, the exact mutually inverse coefficient maps are
\[
\frac{\mathcal T_A(\chi_aP)}{\ell_r}
 =u_r+\sum_{t<r}p_t\frac{\ell_t}{\ell_r}u_t,
\qquad p_t=z_t\frac{\ell_r}{\ell_t}.\tag{CRP8}
\]
The earlier monic minimizers form a triangular basis with diagonal one.
Their conductor images therefore span exactly the preceding relation
space, without dropping any lift. If $\delta_r$ is the degree-order
Cholesky pivot of $H$, its variational definition and CRP8 give
\[
\boxed{\delta_r=\frac{|\ell_r|^2\eta_r}{\nu_r^+},\qquad
\gamma_r=\log\delta_r+\log\nu_r^+-\log\nu^-_{m+r}
                         -2\log|\ell_r|.}\tag{CRP9}
\]
These are full attained pivots. The positive-square minimum is the
classical Schur operation of Boyd and Vandenberghe
\cite[Appendix A.5.5, p.~650]{human:boyd-vandenberghe}; CRP8--9 supplies
its exact original-fibre realization.

\subsection{The killed-space angle from actual root evaluations}
For $v=0$ set $C_{\rm ang}=0$; then $K=I$. Suppose $v>0$ and choose
the following actual upper roots, with no rescaling:
\[
z_h=(2h-1)\gamma+i\delta,\qquad
\beta_h=c+iz_h=\beta_{a,(a-1)/2,(a-1)/2+h},\quad0\le h<v.
\tag{CRP10}
\]
All their indices are in the original grid, since $a\ge257$ and
$v\le80$. Let $\mathcal E=(\phi_{n,s,c}(\beta_h))_{h<v,n\le N}$,
$\mathcal K=\mathcal E\mathcal E^*$, and let $\mathcal L$ be the first
$v$ columns of $\mathcal E$. Evaluation annihilates $\operatorname{im}\Phi$.
The matrix $\mathcal L$ is invertible, because its columns are degree
$0,\ldots,v-1$ polynomials with nonzero leading coefficients at distinct
points. Let $Z$ insert the first $v$ native coordinates. Sylvester's
identity, followed by inclusion of the evaluation adjoint range, gives
\[
\begin{split}
\det K&=\det(I-\Phi^*ZZ^*\Phi)
 =\det(Z^*(I-\Phi\Phi^*)Z),\\
Z^*(I-\Phi\Phi^*)Z&\succeq
 Z^*\mathcal E^*\mathcal K^{-1}\mathcal E Z
 =\mathcal L^*\mathcal K^{-1}\mathcal L,\\
\det K&\ge |\det\mathcal L|^2/\det\mathcal K>0.
\end{split}\tag{CRP11}
\]
For the determinant monotonicity used here, write a positive pair as
$Y^{1/2}(I+C)Y^{1/2}$ with $C\succeq0$ and multiply its positive
eigenvalues. Thus the inequality is also valid when the evaluation
space is a proper part of the full orthogonal complement.

Put
\[
\begin{gathered}
Z_v=\max_{h<v}|z_h|,\quad
\mathcal A=e^2(N+1)^{1+s/2}(2N+1)^{1+Z_v},\\
\mathcal V=(2\gamma)^{v(v-1)/2}\prod_{n=1}^{v-1}n!,\\
C_{\rm ang}=\left[v\log\mathcal A+
 \sum_{n=0}^{v-1}\log(n!(s/2)_n)-2\log\mathcal V\right]_+.
\end{gathered}\tag{CRP12}
\]
Here $\mathcal V$ is the literal Vandermonde modulus of the selected
$z_h$'s. On $|t|=R=N/(N+1)$, the original generating series obeys
\[
|\arctan t|\le\operatorname{artanh}R=\tfrac12\log(2N+1),\qquad
|(1+t^2)^{-s/4}|^2\le(1-R^2)^{-s/2}\le(N+1)^{s/2}.
\]
Cauchy's coefficient formula, $R^{-2n}\le e^2$ for $n\le N$, and
$\rho_n^2\le n!/(1/2)_n\le2n+1$ imply
\[
|\phi_{n,s,c}(\beta_h)|^2\le
 \frac{e^2}{M_s}(N+1)^{s/2}(2N+1)^{1+Z_v}.
\tag{CRP13}
\]
The last bound on $\rho_n$ follows by induction: multiplying $2n+1$
by $(n+1)/(n+1/2)$ gives $2n+2\le2n+3$. Hadamard's inequality is
obtained by orthogonalizing the rows, whose residual lengths do not
exceed their lengths. It gives $\det\mathcal K\le(\mathcal A/M_s)^v$.
The leading coefficients in CRP4 give exactly
$|\det\mathcal L|^2=\mathcal V^2/\prod_{n<v}h_{n,s}$.
Combining them with CRP11 cancels exactly the $v$ factors $M_s$, proving
\[
\boxed{0\le-\log\det K\le C_{\rm ang}
 =O_{\rm actual}(a\log q)=o(q).}\tag{CRP14}
\]
The constant depends on the original fixed quartet and conductor,
including $v$ and $\gamma$. The estimate is uniform in the stated
$M,s$ range; no bound uniform in a varying hypothetical quartet is
claimed.

\subsection{Summed absolute pivots from all exterior ranks}
For clarity the exact finite LT constants used below are
\[
\begin{gathered}
F_0^{\flat}=4+4\sqrt{\delta^2+\gamma^2},\quad
B_0=2^{v+8}\bigl(\sum|a_{uw}|\bigr)e^{2\pi\gamma},\\
B_2=B_0\{F_0^{\flat}(F_0^{\flat}+1)+4vF_0^{\flat}+4v(v+1)\},
\quad H_*=B_0(1+\sqrt3)+2\sqrt{10}B_2,\\
H_s=\begin{cases}\max(1,H_*\sqrt{2v+1}),&s=1,\\
\max(1,H_*),&s=a,\end{cases}\qquad
\kappa_s=\begin{cases}1+\sqrt{8/3},&s=1,\\2,&s=a,\end{cases}\\
J=d\log^+(1/|g_0|)+\tfrac12(s/2-1)_+v\sum_{j=1}^d1/j,\\
F_r=r\log H_s+4(d-r)\log\kappa_s\ (1\le r\le d),
\quad F_0=0,\quad I_r=J+F_{d-r}\ (1\le r\le d),\quad I_0=0.
\end{gathered}\tag{CRP15}
\]
LT7--22/LRC6--9 proves
$\log\|\bigwedge^r A\|\le F_r$ and
$\log\|\bigwedge^r A^{-1}\|\le I_r$ for every original rank.
The same bounds hold for ordinary transposes. Put
$F^*=\max_{0\le r\le M+1}F_r$ and
$I^*=\max_{0\le r\le M+1}I_r$.

We prove the needed selected-pivot estimates, rather than replacing
them by a diagonal trial norm. For an index set $J_+$, the actual
degree-order pivot minimizes the column residual over all earlier
columns. Restricting the minimum to earlier columns in $J_+$ can only
increase it. Multiplying these selected Gram--Schmidt residuals yields
\[
\prod_{r\in J_+}\delta_r\le\det H[J_+,J_+]
 =\|\bigwedge^{|J_+|}(AB_{J_+})\|^2
 \le e^{2F_{|J_+|}}.\tag{CRP16}
\]
The last bound uses that $B_{J_+}$ is a contraction, since the
corresponding columns of $\Phi$ are orthonormal.

For another set $J_-$ first condition every column in $J_-$ on all
columns in its complement, and then on its preceding columns in
$J_-$. These spaces include all original preceding columns, so the
new residuals are no larger. Their product is the Schur determinant
$\det H/\det H[J_-^c,J_-^c]$. The block inverse formula therefore gives
\[
\prod_{r\in J_-}\delta_r^{-1}
 \le\det(H^{-1}[J_-,J_-]).\tag{CRP17}
\]
To estimate the right side, put $B=WK^{1/2}$ where $W^*W=I$, and
$C=W^*A^*AW$. Complete $W$ to an orthonormal basis. If the resulting
blocks of $A^*A$ are $\left(\begin{smallmatrix}C&T\\T^*&R\end{smallmatrix}\right)$,
the original full Schur formula gives
\[
C^{-1}\preceq(C-TR^{-1}T^*)^{-1}=W^*(A^*A)^{-1}W.
\tag{CRP18}
\]
For a zero-dimensional complementary block this is equality. For
any isometric rank-$t$ column matrix $Y$, determinant monotonicity
and exterior norms give
$\det(Y^*C^{-1}Y)\le\|\bigwedge^t A^{-1}\|^2\le e^{2I_t}$.
Taking $Y$ to be a top-$t$ eigenvector frame proves
$\|\bigwedge^t C^{-1}\|\le e^{2I_t}$. Since
$H^{-1}=K^{-1/2}C^{-1}K^{-1/2}$ and every eigenvalue of $K$ lies in
$(0,1]$, the product of any $t$ eigenvalues of $K^{-1}$ is at most
$\det K^{-1}$. Thus
\[
\det(H^{-1}[J_-,J_-])\le e^{2I_{|J_-|}}/\det K.\tag{CRP19}
\]
Choose $J_+=\{r:\log\delta_r>0\}$ and
$J_-=\{r:\log\delta_r<0\}$. Empty sets contribute zero. Taking
logarithms in CRP16--19 and adding proves the finite result
\[
\boxed{\sum_{r=0}^M|\log\delta_r|
 \le 2F^*+2I^*+C_{\rm ang}=O_{\rm actual}(q).}\tag{CRP20}
\]
Indeed $d=O(q)$, $\log H_s,\log\kappa_s=O_{\rm actual}(1)$,
and $J=O_{\rm actual}(q+a\log q)=O_{\rm actual}(q)$.
This controls all pivots together and uses the actual polynomial
ideal position through CRP10--14. It does not estimate the sum of
the trial-to-attained slacks $\log(\tau_r/\eta_r)$ by $O(q)$.

\subsection{Exact directional transfer after the scalar reconstruction}
The companion scalar proof supplies the corrected profile
\[
\begin{split}
c_r&=\Delta\mathfrak f(r/q)
-\tfrac12\{(r+\Delta)\log(1+\Delta/r)-\Delta\},\quad r>0,\\
c_0&=\tfrac\Delta2\{1+\log(q/\Delta)\},\qquad
\mathcal E_c\ge\sum_{r=0}^M|\gamma_r-c_r|
 =O_{\rm actual}(q),\qquad c_r^+=\max(c_r,0).
\end{split}\tag{CRP21}
\]
Here $\mathcal E_c$ is its explicit finite sum of single-norm errors,
CRP20 and scalar correction errors; it is not an assumed smallness
condition. The subsequent proof transfers this proved error without
altering any original directional matrix.

We specify the map into the analytic coefficient frame as well. Let
$e_b=\chi_{a-8}\pi_b^-/\sqrt{\nu_b^-}$ be the original orthonormal
lower-ideal basis. In that basis write the full relation columns as
$\mathcal R_r=[\mathcal L_r;H_r]$, where the first block has $m$ rows.
Degree and the monic coefficient give
$H_r[t,t]=\sqrt{\nu^-_{m+t}}$ and $H_r[t,u]=0$ for $t>u$.
Thus the mutually inverse coefficient maps are $z=H_rb$ and
$b=H_r^{-1}z$, and the relation graph is
$\mathcal R_rH_r^{-1}=[X_r;I]$, $X_r=[x_0,\ldots,x_r]$.
The leading columns persist with $r$ because the inverse of an upper
triangular leading block persists. Put $\Omega_r=I_m+X_rX_r^*$.
For the original analytic columns $F_b=N_b/E_A$ and original
complex-linear functionals $\Lambda_{F_b}(S'^j)=F_b^{(j)}(0)$, define
$A_{\rm an}[h,b]=\Lambda_{F_b}(e_h)$ for $0\le h<m$.
This $m\times m$ matrix is the analytic coefficient map NTG16; the
$d\times d$ matrix $A_{D,s}$ in CRP5 is the conductor on polynomials.
Their exact common intermediate map is the full relation graph just
displayed: CRP8 constructs every relation polynomial before the
graph coordinate change, and functional annihilation gives
$X_r^{\mathsf T}A_{\rm an}+F_{\rm hi}=0$ after it.
NTG18--20's full lower-evaluation minimum therefore gives
$D_r=A_{\rm an}^*\overline{\Omega_r}A_{\rm an}$ in the original
analytic frame, with $D_{-1}=A_{\rm an}^*A_{\rm an}$.
Here the row index notation means exactly
$D_r=D^{\rm nat}_{L+r}=D^{\rm target}_{q+r}$, including $r=-1$.
The analytic frame is onto by the original complete functional map,
so $A_{\rm an}$ is invertible. Moreover
$\det\Omega_r/\det\Omega_{r-1}=\eta_r/\nu^-_{m+r}$: taking the
Gram of $[X_r;I]H_r$ gives
$\det\mathcal R_r^*\mathcal R_r=(\prod_{t\le r}\nu^-_{m+t})\det\Omega_r$,
while completing the same relation square gives
$\det\mathcal R_r^*\mathcal R_r=\prod_{t\le r}\eta_t$.
This proves the morphism carrying the scalar pivots to the native
analytic row, before any directional estimate.

The row after the full original relation minimum is retained with
the ordinary-transpose convention of NTG/NFI:
\[
\mathbf u_r=\frac{\boldsymbol\rho_{L+r}}{\sqrt{1+\zeta_{L+r}}}
   =-i^{-(L+r)}x_r^{\mathsf T}A_{\rm an},\quad
D_r=D_{r-1}+\mathbf u_r^*\mathbf u_r,\quad
t_r=\mathbf u_rD_{r-1}^{-1}\mathbf u_r^*=e^{\gamma_r}-1.\tag{CRP22}
\]
The symbol $\boldsymbol\rho_{L+r}$ here is the original residual \emph{row};
the scalar $\rho_n$ in CRP4--5 has already served only to specify the
conductor matrix. Their original distinction is preserved.
The full relation Gram defining $\zeta_{L+r}$ and this row is the
one used in ATA8 and NFI, not its diagonal. For the four-copy row
$U_r=\operatorname{diag}(\mathbf u_r,4)$ let
$\mathcal D_0=\operatorname{diag}(D_{r-1},4)$.
For any actual flag column map $X$, the determinant update is
\[
\begin{split}
K_{X,r}&=U_rX(X^*\mathcal D_0X)^{-1}X^*U_r^*,\\
\Delta\log\det(X^*\mathcal D X)&=\log\det(I_4+K_{X,r}).
\end{split}\tag{CRP23}
\]
The determinant identity follows by factoring the old positive Gram
and applying $\det(I+AB)=\det(I+BA)$; the inverse rank-one update is
the classical formula recorded by Hager \cite[p.~221, (2)]{human:hager1989}.
For empty flags the matrices and increments are zero.

If $t_r>0$, define
$Z_r=t_r^{-1/2}U_r\mathcal D_0^{-1/2}$ and let $P_{X,r}$ be the
orthogonal projector onto $\mathcal D_0^{1/2}\operatorname{im}X$.
Then
\[
Z_rZ_r^*=I_4,\quad M_{X,r}=Z_rP_{X,r}Z_r^*,\quad
0\preceq M_{X,r}\preceq I_4,\quad K_{X,r}=t_rM_{X,r}.
\tag{CRP24}
\]
All four eigenvalues $\lambda_{X,r,b}$ of $M_{X,r}$, including zero,
are retained. If $t_r=0$, the original row is zero, and all its exact
and approximating directional contributions are defined to be zero;
no direction is assigned to a zero row.

For $g\ge0$ define $F(g,\lambda)=\log(1+(e^g-1)\lambda)$ and
$C(g,\lambda)=(g+\log\lambda)_+$, with $C(g,0)=0$.
Direct differentiation and the two cases $e^g\lambda\le1$ and
$e^g\lambda\ge1$ give
\[
0\le F-C\le\log2,\qquad
\partial_gF=\frac{e^g\lambda}{1+(e^g-1)\lambda}\in[0,1].
\tag{CRP25}
\]
Moreover $\partial_gF$ is increasing in $\lambda$. The difference
$d_g(\lambda)=F-C$ increases from zero up to
$\lambda=e^{-g}$ and decreases thereafter to zero at $\lambda=1$.
At $g=0$ it vanishes identically. Hence it can be written as
$d_g=u_g-v_g$, with $u_g,v_g$ increasing, both starting at zero and
both ending at the same maximum $\log(2-e^{-g})\le\log2$.

The original two chains are
\[
J\subset W_i\subset W\subset T,\qquad
K_i\subset I\subset K_c\subset S.\tag{CRP26}
\]
Projection order and the Rayleigh min--max principle imply order of
the four eigenvalues at each position along each chain. For four
ordered scalar arguments $z_1\le z_2\le z_3\le z_4$, the sum of
the two disjoint increments of any increasing $u$ is between zero
and $u(1)-u(0)$. Applying this first to $\partial_gF$ shows that the
derivative of each chain sum is in $[0,4]$. Their difference,
with the exact signs below, is therefore $4$-Lipschitz in $g$.
Applying it to $u_g$ and $v_g$ shows that the discrepancy of each
chain sum lies between $-4\log2$ and $4\log2$. Thus the difference
of the two chains has discrepancy at most $8\log2$.

Retain the eight \emph{labelled occurrences}, even if some underlying
subspaces coincide:
\[
(T,K_c,W_i,K_i,W,S,J,I),\qquad
(\sigma_X)=(+,+,+,+,-,-,-,-).\tag{CRP27}
\]
For any original row set $\mathscr I$ with its original weights
$0\le\vartheta_r\le2$, put
$\widetilde f_{X,r}=\sum_{b=1}^4C(c_r^+,\lambda_{X,r,b})$ on nonzero
rows. Since $\gamma_r\ge0$,
$|\gamma_r-c_r^+|\le|\gamma_r-c_r|$. Freeze the actual matrices
$M_{X,r}$ while applying the scalar Lipschitz estimate; this does
not assert that the flags of a different metric equal these flags.
Summation proves
\[
\boxed{\left|\mathcal L_e^{\mathscr I}-
 \sum_{r\in\mathscr I}\vartheta_r\sum_X\sigma_X\widetilde f_{X,r}\right|
 \le8\mathcal E_c+8\log2\sum_{r\in\mathscr I}\vartheta_r.}
\tag{CRP28}
\]
For a single original nested quotient pair $A\subset B$, the
derivative of its four-eigenvalue difference is in $[0,4]$ and its
clipped discrepancy is at most $4\log2$. Therefore
\[
\boxed{|V_{B/A}^{\mathscr I}-\widetilde V_{B/A}^{\mathscr I}|
 \le8\mathcal E_c+4\log2\sum_{r\in\mathscr I}\vartheta_r.}
\tag{CRP29}
\]
This applies to STF's unchanged complementary pairs $W/W_i$,
$U/T$, $S/K_c$, and $I/K_i$, where $U$ is the full original ambient
space and its directional compression is $I_4$. It includes the
original quotient minima through the determinant differences in
CRP23, not new quotient ranks or substituted forms.

\subsection{One-copy allocations, exact quotient maps, and backward use}
Let $B_a$ be the original onto observation map, $J_a$ an injective
matrix with image $\ker B_a$, and let the fixed onto-value frame be
unchanged. Its restricted and quotient forms are exactly
\[
G_{R,r}=J_a^*D_rJ_a,\qquad
G_{K,r}=(B_aD_r^{-1}B_a^*)^{-1}.\tag{CRP30}
\]
If $Y$ is any right inverse of $B_a$, the fixed square frame
$[J_a,Y]$ has determinant independent of $r$. Completing the full
$J_a$ coefficient square proves
\[
|\det[J_a,Y]|^2\det D_r=\det G_{R,r}\det G_{K,r}.
\tag{CRP31}
\]
Thus its original four-endpoint return loses no frame factor.
For a nonzero row let $P_{R,r}$ be the orthogonal projector onto
$D_{r-1}^{1/2}\ker B_a$ and retain the actual projection energy
$p_r=t_r^{-1}\mathbf u_rD_{r-1}^{-1/2}P_{R,r}
D_{r-1}^{-1/2}\mathbf u_r^*\in[0,1]$.
Then CRP23 in one copy and CRP31 give
\[
\Delta S_R=F(\gamma_r,p_r),\qquad
\Delta S_K=\gamma_r-F(\gamma_r,p_r).\tag{CRP32}
\]
Set $\kappa_r=-\log p_r$, allowing $\kappa_r=+\infty$ when $p_r=0$.
Both functions in CRP32 are $1$-Lipschitz in the scalar argument.
CRP25 and $g-C(g,p)=\min(g,-\log p)$ now prove
\[
\begin{gathered}
\widetilde S_R=\sum_r\vartheta_r(c_r^+-\kappa_r)_+,\qquad
\widetilde S_K=\sum_r\vartheta_r\min(c_r^+,\kappa_r),\\
\boxed{|S_X-\widetilde S_X|\le2\mathcal E_c+
 \log2\sum_r\vartheta_r,\qquad X=K,R.}\tag{CRP33}
\end{gathered}
\]
Zero rows contribute zero in both expressions, as above. The exact
identities remain $S_K+S_R=\mathcal T$ and
$4S_K=\mathcal L_e+V_b+V_{\rm rem}+V_s+V_A$; every term uses the same
original row and quotient maps. CRP28--33 do not evaluate the
$p_r$ or four-dimensional spectra from dimensions or from the total.

The precise replacement in older PCB/LRS/FTR row receivers is the
following: replace the radial function $\Delta\mathfrak f(r/q)$
by CRP21's $c_r$ (and by $c_r^+$ only inside clipped directional
expressions), replace its weighted $O(q\log q)$ error by the
explicit $2\mathcal E_c$, and use CRP28--33 with the original
weights, flags and quotient pairs. Keeping the old uncorrected
profile and merely changing its allowance to $O(q)$ is invalid:
the subtracted correction has full-window weighted sum
$128q\log a+O(q)$, proved in the companion scalar calculation.
For prefix/tail receivers also replace the integrated centres by
the companion $\mathcal A_\Delta$ formulas, including their endpoint
terms. The old trial-to-attained and metric-chain comparisons remain
their own proved comparisons. Every old sealed edition is preserved;
this is a successor theorem with an explicit map back to each receiver.
\endgroup
